\chapter{The Caster's Ledger}

A spell you cast once is a demonstration. A spell you cast twice is a habit. A spell you record is a \emph{tool}. The Weave does not care whether you write things down, but you should. Memory is a poor ledger. Paper is patient.

This appendix provides a template for building your own grimoire---a personal record of the spells you have learned, invented, or stolen. I have left space for the spell's name, its tags, its intended effect, and a note on what the Weave said in reply. The last column is the most important.

\section{How to Use This Ledger}

Before you cast a new spell, write down your intention. Name the tags. Estimate the Difficulty. This is not a contract; it is a \emph{hypothesis}.

After you cast, record what happened. Did the Weave comply? Did it comment? Did it laugh? Write that down too.

Over time, you will notice patterns. Some tags work well together. Some elements produce reliable backlash. Some spells fail in instructive ways. Your ledger will become a map of your relationship with the Weave.

I have kept such a ledger for forty years. It is thicker than my arm. I do not carry it; it carries me.

\section{Spell Card Template}

The following page contains a template for a spell card. You may photocopy it, trace it, or simply write its contents in a notebook. I have provided a sample entry to illustrate the form.

\subsection{Sample Entry}

\noindent
\begin{tabular}{|p{0.45\textwidth}|p{0.45\textwidth}|}
\hline
\multicolumn{2}{|c|}{\textbf{Spell Name: Ember Flick}}\\
\hline
\textbf{Tags} & [BURNING] [STRIKE] \\
\hline
\textbf{Intended Effect} & Ignite a small flammable object within Close range. A candle, a twist of paper, a dry leaf. Not a person. \\
\hline
\textbf{Difficulty (estimated)} & 3 (two tags, no dangerous surcharge) \\
\hline
\textbf{Outcome \& Backlash} & Clean success: the leaf ignites. Minor backlash on a partial: a spark jumps to the caster's sleeve (1 Fatigue). \\
\hline
\textbf{Notes} & The Weave is precise about what it considers ``flammable.'' Wet leaves do not count. Do not argue. \\
\hline
\end{tabular}

\subsection{Blank Spell Card}

Reproduce this card as many times as you need. I recommend a stout paper and a pencil with a good eraser. The Weave appreciates correction.

\noindent
\begin{tabular}{|p{0.45\textwidth}|p{0.45\textwidth}|}
\hline
\multicolumn{2}{|c|}{\textbf{Spell Name: }}\\
\hline
\textbf{Tags} & \\
\hline
\textbf{Intended Effect} & \\
\hline
\textbf{Difficulty (estimated)} & \\
\hline
\textbf{Outcome \& Backlash} & \\
\hline
\textbf{Notes} & \\
\hline
\end{tabular}

\section{Advanced Ledger: The Spell Chain}

Some spells are not single castings. They are sequences---a [VEIL] followed by a [LEAP] followed by a [SILENCE]. The Weave treats each step as a separate transaction, but you may record the chain as a single entry.

\subsection{Sample Chain: The Quiet Exit}

\noindent
\begin{tabular}{|p{0.45\textwidth}|p{0.45\textwidth}|}
\hline
\multicolumn{2}{|c|}{\textbf{Chain Name: The Quiet Exit}}\\
\hline
\textbf{Step 1} & [VEIL] on self. DV 2. Blur features. \\
\textbf{Step 2} & [SILENCE] on footsteps. DV 2. Muffle sound. \\
\textbf{Step 3} & [LEAP] to exit. DV 4 (dangerous). Teleport to visible exit. \\
\hline
\textbf{Total Risk} & Each step carries its own backlash. A failure in Step 3 does not undo Step 1. The Weave respects sequence. \\
\hline
\textbf{Notes} & Practice the chain in a safe room before attempting it in a guarded hallway. I speak from experience. \\
\hline
\end{tabular}

\section{The Ledger's Discipline}

A student once showed me his spell ledger. It was beautiful---leather binding, gold leaf, neat calligraphy. He had recorded every spell he had ever cast, including the ones that failed.

I asked him, \emph{``Do you read it?''}

He said, \emph{``I wrote it. I do not need to read it.''}

He was wrong. A ledger that is not consulted is a monument, not a tool. A ledger that is consulted is a conversation with your past self. The Weave listens when you learn from your own mistakes.

\subsection{Monthly Review}

Once a month, read your ledger from cover to cover. Note which spells succeed reliably. Note which spells fail in interesting ways. Note which spells you have not cast since you wrote them down---and ask yourself why.

\subsection{The Weave's Marginalia}

Sometimes the Weave adds its own notes. I have opened my ledger to find a word underlined in a hand that was not mine. I have seen a page that was blank become filled with a single sentence: \emph{``You forgot the [WITNESS].''}

The Weave does not vandalise. It \emph{edits}. Welcome the edits.

\section{A Final Card for the Road}

Carry one blank card with you at all times. When you cast a spell you have never cast before, write it down immediately---the tags, the intent, the backlash. Do this while the Weave's receipt is still warm.

The card will become a record of your most desperate moments. Those are the moments that teach the most.

\vspace{2em}

\noindent
\begin{tabular}{|p{0.45\textwidth}|p{0.45\textwidth}|}
\hline
\multicolumn{2}{|c|}{\textbf{Emergency Spell}}\\
\hline
\textbf{Tags} & \\
\hline
\textbf{Intended Effect} & \\
\hline
\textbf{Difficulty} & \\
\hline
\textbf{Backlash Observed} & \\
\hline
\textbf{Date} & \\
\hline
\end{tabular}

\vspace{1em}

\textit{--- Lysandra}
